\documentclass[10pt,a4paper]{amsart}
\usepackage[margin=19mm]{geometry}
\usepackage{amsmath,amssymb,mathtools}
\usepackage[scr=rsfs,cal=euler]{mathalfa}
\usepackage{xfrac}
\usepackage[unicode,hidelinks,bookmarks=false]{hyperref}
\newcommand{\ie}{i.e.\ }
\newcommand{\cf}{cf.\ }
\newcommand{\resp}{respectively\ }
\newcommand{\Subsection}{\subsection}
\providecommand{\nobreakdash}{\nobreak\hbox{-}}
\setlength{\emergencystretch}{2em}
\multlinegap0pt

\theoremstyle{plain}
\newtheorem{theorem}{Theorem}[section]
\newtheorem{lemma}[theorem]{Lemma}
\newtheorem{prop}[theorem]{Proposition}
\newtheorem{corollary}[theorem]{Corollary}
\newtheorem{assumption}[theorem]{Assumption}
\newtheorem{definition}[theorem]{Definition}
\newtheorem{remark}[theorem]{Remark}
\newtheorem{example}[theorem]{Example}
\title[Global bifurcation of water-wave fronts generated by surface pressure]{Global bifurcation of water-wave fronts generated by surface pressure}
\author{Jifeng Chu, Zihao Wang, Yong Zhang}
\date{Source: arXiv:2609.28024v1 (CC BY 4.0)}
\begin{document}
\maketitle

\noindent\emph{Source and licence.} The delivered work is the article shown in the title above, version arXiv:2609.28024v1 (CC BY 4.0), distributed under the Creative Commons Attribution 4.0 International licence, as recorded in the arXiv licence metadata for this exact version and shown on the abstract page saved with this item. The full licence notice, which carries the licence URL and the address of that abstract page, is the bundled file \texttt{licenses/arxiv-2609.28024-CC-BY-4.0.txt}.
\par\smallskip

\noindent\textit{Editorial scope.} Counted unit: the article's Lemma (source label lem:range-regularity) (source0.tex lines 906--914, source label \texttt{lem:range-regularity}): ``Assume , and let denote the interior and boundary data. If satisfies , then . Consequently, L( X ) L( X b) Y , ( L X b ) ( L X ).''. It is delivered together with the article's complete original proof (source0.tex lines 915--999, one proof environment adjacent to the statement, 85 lines), copied line for line from the article's own text. Required context retained uncounted: eq:two-ended-space (lines 390--397); lem:limit-invertibility (lines 782--786). Editorial formatting only: an AMS \texttt{amsart} preamble that restates the article's own macro definitions and its own theorem declarations, and the theorem environments the retained lines use; the article's own class file is neither shipped nor input; the retained environments are renumbered sequentially in order of appearance, whereas the article prints them with its own numbering; the article's equation numbering is not reproduced and every cross-reference inside the retained text resolves to the wrapper's numbering. No citation occurs inside any retained line, so the wrapper carries no bibliography; All retained bodies are exact copies of the bundled \texttt{sources/211545/source0.tex}, so no omission receipt applies. No asset-inclusion command occurs inside any retained span and the package ships no image asset.


\section*{Retained context: eq:two-ended-space (source0.tex lines 390--397)}

	\begin{equation}\label{eq:two-ended-space}
		\begin{aligned}
			C_\infty^{k+\alpha}(\overline\Omega):=
			\bigl\{\chi_-u_-+\chi_+u_++u_0:{}&
			u_\pm\in C^{k+\alpha}([-m,0]),~u_0\in C_b^{k+\alpha}(\overline\Omega)
			\cap C_0^k(\overline\Omega)\bigr\}.
		\end{aligned}
	\end{equation}


\section*{Retained context: lem:limit-invertibility (source0.tex lines 782--786)}

	\begin{lemma}\label{lem:limit-invertibility}
	If $\sigma_0^{\pm}<0$, then
	$\mathscr{L}_{\pm}:\mathscr{X}_b\to \mathscr{Y}_b$
	is an isomorphism.
	\end{lemma}


\section{The counted unit: Lemma (source label lem:range-regularity) and its complete original proof}

	\begin{lemma}\label{lem:range-regularity}
	{Assume $\sigma_0^\pm<0$, and let $f=(f_1,f_2)\in\mathscr Y_\infty$ denote the interior and boundary data. If $\phi\in\mathscr X_b$ satisfies $\mathscr L\phi=f$, then $\phi\in\mathscr X_\infty$. Consequently,}
	\begin{equation}\label{eq:range-intersection}
		\mathscr L(\mathscr X_\infty)=\mathscr L(\mathscr X_b)\cap\mathscr Y_\infty,
		\qquad
		\ker(\mathscr L|_{\mathscr X_b})
		=\ker(\mathscr L|_{\mathscr X_\infty}).
		\end{equation}
	\end{lemma}

	\begin{proof}
		{We first prove convergence at $+\infty$ for $\phi$ and its derivatives of order at most three.} We claim that for any multi-index $\nu=(\nu_q,\nu_p)$ with $|\nu|\leq3$, the family $\{D^\nu \phi(q,\cdot):q\geq Q\}$ is Cauchy in $C^0([-m,0])$ as $Q\to\infty$. If not, there exist some $\epsilon>0$, sequences $q_{1n},q_{2n}\to\infty$ and $p_n\to p_*\in[-m,0]$ such that
		\begin{equation}\label{eq:endpoint-separation}
		|D^\nu\phi(q_{1n},p_n)-D^\nu\phi(q_{2n},p_n)|\geq\epsilon.
		\end{equation}
		{Define the two translated functions $\phi_n^{(k)}$ and their difference $v_n$ by}
		\begin{equation*}
			\phi_n^{(k)}=\phi(q+q_{kn},p),\qquad v_n=\phi_n^{(1)}-\phi_n^{(2)},\quad k=1,2,
		\end{equation*}
		{Let $\mathscr L_n^{(k)}$ be the operator obtained by translating the coefficients of $\mathscr L$ by $q_{kn}$. Subtracting the two translated equations, including their boundary components, gives}
		\begin{equation}\label{eq:translated-difference-equation}
			{\mathscr L_n^{(1)}}v_n
			=f(\,\cdot+q_{1n},\cdot)-f(\,\cdot+q_{2n},\cdot)
			+({\mathscr L_n^{(2)}}-{\mathscr L_n^{(1)}})\phi_n^{(2)}.
		\end{equation}
		Because \(h\in H+\mathscr X_\infty\), the coefficients of both
		\({\mathscr L_n^{(i)}}\), together with the derivatives needed to pass to the
		divergence equation and oblique boundary condition, converge uniformly on
		compact subsets to those of \(\mathscr L_+\).  Since \(f\in\mathscr Y_\infty\), the
		first difference on the right of
		\eqref{eq:translated-difference-equation} tends to zero locally, together
		with the corresponding interior and boundary derivatives.  The second
		difference also tends to zero locally: the coefficient difference converges
		to zero, while \(\phi_n^{(2)}\) is uniformly bounded in
		\(C_b^{3+\alpha}\).
		
		The sequence \(v_n\) is bounded in \(C_b^{3+\alpha}\).  Local compact 	embedding, including the Dirichlet and oblique boundaries, therefore gives a	subsequence such that
		\begin{equation*}
			v_n\longrightarrow v\quad\hbox{in }C_{\rm loc}^{3}(\overline\Omega),
			\qquad {\mathscr L_+}v=0,\qquad v=0\ \hbox{on }\mathcal B .
		\end{equation*}
		The uniform H\"older bounds pass to the limit, so \(v\in\mathscr X_b\); in particular, it is bounded on the whole strip.  Lemma~\ref{lem:limit-invertibility}
		forces \(v=0\), whereas \eqref{eq:endpoint-separation} and the
		\(C^3_{\rm loc}\) convergence give \(|D^\nu v(0,p_*)|\ge\epsilon\).
		This contradiction proves our claim for every
		\(|\nu|\le3\).
		
		Let \(\ell_\nu(p)\) be the uniform limit of $D^\nu\phi(q,p)$.  If \(\nu_q\ge1\), then
		\(D^{\nu-(1,0)}\phi\) is bounded and its \(q\)-derivative tends uniformly to
		\(\ell_\nu\).  A nonzero value of \(\ell_\nu(p)\) would make
		\(D^{\nu-(1,0)}\phi(q,p)\) unbounded, after fixing
		\(p\) and integrating with respect to \(q\).   Hence
		\begin{equation}\label{eq:tangential-limits-zero}
			\ell_\nu\equiv0\qquad\text{whenever }\nu_q\ge1,\quad |\nu|\le3.
		\end{equation}
		Define \(\phi_+(p):=\ell_{(0,0)}(p)\).  Uniform convergence of the
		{$p$-derivatives and the fundamental theorem of calculus show that}
		\[
		\phi_+\in C^3([-m,0]),\qquad
		\partial_p^j\phi_+=\ell_{(0,j)},\quad 0\le j\le3.
		\]
		The uniform \(C_b^{3+\alpha}\) bound for \(\phi\) passes to the limit in
		the H\"older seminorm, so in fact \(\phi_+\in C^{3+\alpha}([-m,0])\).
	 {Let $f_{1,+}$ and $f_{2,+}$ denote the downstream limits of the interior and boundary data. Passing to the limit in the equations yields}
		\begin{equation*}
			(K_{p,+}^{-3}\phi_{+,p})_p=f_{1,+},\qquad
			-K_{p,+}(0)^{-3}\phi_{+,p}(0)+g\phi_+(0)=f_{2,+},\qquad
			\phi_+(-m)=0 .
		\end{equation*}
	{The Cauchy property proved above and \eqref{eq:tangential-limits-zero} give}
		\begin{equation}\label{eq:full-endpoint-convergence}
			\max_{i+j\le3}\sup_{-m\le p\le0}
			\left|\partial_q^i\partial_p^j
			\bigl(\phi(q,p)-\phi_+(p)\bigr)\right|
			\longrightarrow0\qquad(q\to+\infty).
		\end{equation}
		
		The same argument at \(-\infty\) produces
		\(\phi_-\in C^{3+\alpha}([-m,0])\) and the analogue of
		\eqref{eq:full-endpoint-convergence}. With the fixed cutoffs from
		{\eqref{eq:two-ended-space}, define the decaying remainder $\phi_0$ by}
		\[
		\phi_0:=\phi-\chi_-\phi_- -\chi_+\phi_+.
		\]
		Then \(\phi_0\in C_b^{3+\alpha}(\overline\Omega)\), and every derivative of
		\(\phi_0\) of order at most three tends uniformly to zero at both ends.
		Thus \(\phi_0\in C_b^{3+\alpha}\cap C_0^3\), proving
		\(\phi\in\mathscr X_\infty\).
		Finally,
		\(\mathscr L(\mathscr X_\infty)\subset\mathscr L(\mathscr X_b)\cap\mathscr Y_\infty\) is
		immediate.  Conversely, if an element of the intersection equals
		\(\mathscr L\phi\) for some \(\phi\in\mathscr X_b\), the result just proved puts
		\(\phi\) in \(\mathscr X_\infty\).  Taking the right-hand side to be zero gives the
		kernel identity in \eqref{eq:range-intersection}.	
	\end{proof}

\end{document}
